% Part: lambda-calculus
% Chapter: syntax
% Section: alpha

\documentclass[../../../include/open-logic-section]{subfiles}

\begin{document}

\olfileid{lam}{syn}{alp}
\olsection{تحويل~$\alpha$}

لننظر في العلاقة بين~$\lambd[x][x]$ و$\lambd[y][y]$: كلاهما يمثل دالة الهوية،
ولكنهما حدان متغايران من حيث التركيب. وليس بينهما اختلاف إلا اسم المتغير المقيد؛
فأحدهما يحصل من الآخر بتغيير هذا الاسم. وهذا هو \emph{تحويل~$\alpha$}.

\begin{defn}[تغيير المتغير المقيد، $\aconvone$]
  إذا احتوى الحد~$M$ على واقعة لـ$\lambd[x][N]$، وتحقق~$y \notin
  \FV{N}$، وكان~$\Subst{N}{y}{x}$ معرّفًا، ووضعنا موضع تلك الواقعة
  \begin{equation*}
    \lambd[y][\Subst{N}{y}{x}]
  \end{equation*}
  فحصل الحد~$M'$، سمينا العملية \emph{تغييرًا للمتغير المقيد}، ورمزنا إليها بـ
  $M \redone[\alpha] M'$.
\end{defn}

\begin{defn}[توافق العلاقة]
  العلاقة~$R$ على الحدود \emph{متوافقة} متى تحققت فيها الشروط الثلاثة:
  \begin{enumerate}
  \item من~$R N N'$ يلزم~$R \lambd[x][N] \lambd[x][N']$
  \item من~$R P P'$ يلزم~$R (PQ) (P'Q)$
  \item من~$R Q Q'$ يلزم~$R (PQ) (PQ')$
  \end{enumerate}
\end{defn}

وبهذا نصوغ التعريف على وجه آخر:
\begin{defn}[تغيير المتغير المقيد، $\aconvone$]
  علاقة \emph{تغيير المتغير المقيد} ($\redone[\alpha]$) أصغر العلاقات المتوافقة
  على الحدود التي تستوفي الشرط التالي:
  \begin{align*}
    &\lambd[x][N] \redone[\alpha] \lambd[y][\Subst{N}{y}{x}] && \text{إذا
      كان~$x \neq y$ و$y \notin \FV{N}$} \\
    & &&\text{وكان~$\Subst{N}{y}{x}$ معرّفًا}
  \end{align*}
\end{defn}

ومعنى «أصغر» أنها تضم ما يوجبه التوافق والشرط المذكور من الأزواج، ولا تضم ما زاد عليه.
ومن هنا نحصل على التعريف الآتي أيضًا:

\begin{defn}[تغيير المتغير المقيد، $\aconvone$] \ollabel{defn:aconvone}
  لعلاقة \emph{تغيير المتغير المقيد} ($\aconvone$) القواعد الاستقرائية التالية:
  \begin{enumerate}
  \item من~$N \aconvone N'$ يلزم~$\lambd[x][N] \aconvone
    \lambd[x][N']$ \ollabel{defn:aconvone1}
  \item من~$P \aconvone P'$ يلزم~$(PQ) \aconvone (P'Q)$ \ollabel{defn:aconvone2}
  \item من~$Q \aconvone Q'$ يلزم~$(PQ) \aconvone (PQ')$ \ollabel{defn:aconvone3}
  \item إذا تحقق~$x \neq y$ و$y \notin \FV{N}$ وكان~$\Subst{N}{y}{x}$
    معرّفًا، لزم~$\lambd[x][N] \redone[\alpha]
    \lambd[y][\Subst{N}{y}{x}]$.
    \ollabel{defn:aconvone4}
  \end{enumerate}
\end{defn}

هذه التعريفات متكافئة، وإثبات تكافؤها متروك تمرينًا. والذي نعتمده فيما يأتي هو التعريف الاستقرائي.

\begin{defn}[تحويل~$\alpha$، $\aconv$]
  علاقة \emph{تحويل~$\alpha$} ($\aconv$) أصغر علاقة انعكاسية متعدية على الحدود
  تشتمل على~$\aconvone$.
\end{defn}

والمراد بـ«أصغر»، على ما تقدم، ألا تضم العلاقة إلا ما يوجبه الانعكاس والتعدي
واحتواء~$\aconvone$. وهذا يعطي التعريف المكافئ التالي:

\begin{defn}[تحويل~$\alpha$، $\aconv$] \ollabel{defn:aconv}
  لعلاقة \emph{تحويل~$\alpha$} ($\aconv$) التعريف الاستقرائي الآتي:
  \begin{enumerate}
  \item من~$P \aconv Q$ و$Q \aconv R$ يلزم~$P \aconv R$.
    \ollabel{defn:aconv1}
  \item من~$P \aconvone Q$ يلزم~$P \aconv Q$. \ollabel{defn:aconv2}
  \item $P \aconv P$. \ollabel{defn:aconv3}
  \end{enumerate}
\end{defn}

\begin{ex}
  الحد~$\lambd[x][f x]$ يقبل تحويل~$\alpha$ إلى~$\lambd[y][f
  y]$، وكذلك العكس.
  والمعنى غير الصوري أن كليهما دالة تأخذ وسيطًا وتردّ حاصل تطبيق~$f$ عليه،
  مع الإحالة إلى متغير البيئة~$f$.

  وأما~$\lambd[x][f x]$ فلا يقبل تحويل~$\alpha$ إلى~$\lambd[x][g
    x]$؛
  لأنهما، بالمعنى غير الصوري، يحيلان إلى متغيري البيئة~$f$ و$g$ بالترتيب.
  فهما دالتان مختلفتان: إذا اختلف~$f$ و$g$ في البيئة اختلف سلوكهما فيها.
\end{ex}

\begin{prob}
  بيّن في كل زوج مما يلي هل يقبل حداه التحويل بحسب~$\alpha$:
  \begin{enumerate}
  \item $\lambd[x][\lambd[y][x]]$ و$\lambd[y][\lambd[x][y]]$
  \item $\lambd[x][\lambd[y][x]]$ و$\lambd[c][\lambd[b][a]]$
  \item $\lambd[x][\lambd[y][x]]$ و$\lambd[c][\lambd[b][a]]$
  \end{enumerate}
\end{prob}

\begin{lem}\ollabel{lem:fv-one}
  من~$P \aconvone Q$ يلزم~$\FV{P} = \FV{Q}$.
\end{lem}

\begin{proof}
  نستقرئ~!!{derivation} العبارة~$P \aconvone Q$.
  \begin{enumerate}
  \item إذا ختم الاشتقاق بالقاعدة~\olref{defn:aconvone4}، كان~$P$ على
    صورة~$\lambd[x][N]$ و$Q$ على صورة
    $\lambd[y][\Subst{N}{y}{x}]$، حيث~$x \neq y$ و$y \notin \FV{N}$
    وكان~$\Subst{N}{y}{x}$ معرّفًا. ثم نفصل بحسب~$x \in \FV{N}$:
    \begin{enumerate}
    \item إن كان~$x \in \FV{N}$، حسبنا
      \begin{align*}
        \FV{\lambd[y][\Subst{N}{y}{x}]} &=
        \FV{\Subst{N}{y}{x}} \setminus \{y\} \\
        & = ((\FV{N} \setminus \{x\}) \cup \{y\}) \setminus \{y\}
         && \text{ بموجب \olref[sub]{thm:infv}} \\
        & = \FV{N} \setminus \{x\} \\
        & = \FV{\lambd[x][N]}
      \end{align*}
    \item وإن كان~$x \notin \FV{N}$، حسبنا
      \begin{align*}
        \FV{\lambd[y][\Subst{N}{y}{x}]}
        & = \FV{\Subst{N}{y}{x}} \setminus \{y\} \\
        & = \FV{N} \setminus \{x\}
         && \text{بموجب \olref[sub]{thm:notinfv}} \\
        & = \FV{\lambd[x][N]}.
      \end{align*}
    \end{enumerate}
  \item وأما الحالات الثلاث الباقية فإثباتها تمرين.
  \end{enumerate}
\end{proof}

\begin{prob}
  أتمّ الحالات المتروكة من برهان~\olref[lam][syn][alp]{lem:fv-one}.
\end{prob}

\begin{lem}\ollabel{lem:inv}
  من~$P \aconvone Q$ يلزم~$Q \aconvone P$.
\end{lem}

\begin{proof}
  نستقرئ~!!{derivation} العبارة~$P \aconvone Q$.
  \begin{enumerate}
  \item إذا ختم الاشتقاق بالقاعدة~\olref{defn:aconvone4}، كان~$P$ على
    صورة~$\lambd[x][N]$ و$Q$ على صورة
    $\lambd[y][\Subst{N}{y}{x}]$، حيث~$x \neq y$ و$y \notin \FV{N}$
    وكان~$\Subst{N}{y}{x}$ معرّفًا. فنجد أولًا~$x \notin
    \FV{\Subst{N}{y}{x}}$ بموجب~\olref[sub]{thm:clr}.
    ثم تقتضي \olref[sub]{thm:inv} أن~$\Subst{\Subst{N}{y}{x}}{x}{y}$
    ليس معرّفًا فحسب، بل يساوي~$N$ أيضًا. وبالقاعدة
    \olref{defn:aconvone4} ينتج~$\lambd[y][\Subst{N}{y}{x}]
    \aconvone \lambd[x][\Subst{\Subst{N}{y}{x}}{x}{y}] =
    \lambd[x][N]$.
    % تصحيح مصدر (0362-I1): المتغير الذي تزيله لمة الخلو من الإحلال الناتج هو x، وهو المطلوب لخطوة ألفا العكسية.
  \item وإذا كانت القاعدة الأخيرة إحدى قواعد التوافق الثلاث، طبقنا فرض
    الاستقراء على التحويل الواقع في متن التجريد أو في الطرف الموافق من
    التطبيق، ثم طبقنا قاعدة التوافق نفسها في الاتجاه المعاكس.
  \end{enumerate}
\end{proof}

\begin{prob}
  أتمّ برهان~\olref[lam][syn][alp]{lem:inv}.
\end{prob}

\begin{thm}
  علاقة تحويل~$\alpha$ تكافؤ على الحدود: فهي انعكاسية، متناظرة، متعدية.
\end{thm}

\begin{proof}
  \begin{enumerate}
  \item كل حد~$M$ يرجع إلى~$M$ من غير تغيير للمتغيرات المقيدة،
    أي بإجراء \emph{صفر} من التغييرات.
  \item إذا انتقل~$P$ بتحويل~$\alpha$ إلى~$Q$ بسلسلة من تغييرات المتغير المقيد،
    بدأنا من~$Q$ وعكسنا كل تغيير بموجب~\olref{lem:inv}،
    وأجرينا المعكوسات بترتيب معاكس، فعُدنا إلى~$P$.
  \item إذا انتقل~$P$ بتحويل~$\alpha$ إلى~$Q$ بسلسلة،
    ثم انتقل~$Q$ إلى~$R$ بسلسلة أخرى،
    انتقل~$P$ إلى~$R$ بإجراء السلسلة الأولى ثم الثانية.
  \end{enumerate}
\end{proof}

ومن الآن نسمي~$M$ و$N$ \emph{متكافئين بحسب~$\alpha$}، ونكتب
$M \aeq N$، متى أمكن تحويل~$M$ بحسب~$\alpha$ إلى~$N$.
وقد أثبتنا أن ذلك يكافئ إمكان تحويل~$N$ بحسب~$\alpha$ إلى~$M$.

\begin{thm}\ollabel{thm:fv}
  من~$M \aeq N$ يلزم~$\FV{M} = \FV{N}$.
\end{thm}

\begin{proof}
  هذه نتيجة مباشرة لـ\olref{lem:fv-one}.
\end{proof}

\begin{lem}\ollabel{lem:sub:R}
  إذا تحقق~$R \aeq R'$ وكان~$\Subst{M}{R}{y}$ معرّفًا، كان
  $\Subst{M}{R'}{y}$ معرّفًا كذلك، وكان مكافئًا بحسب~$\alpha$ للحد~$\Subst{M}{R}{y}$.
\end{lem}

\begin{proof}
  يترك إثباتها تمرينًا.
\end{proof}

\begin{prob}
  أقم برهان~\olref[lam][syn][alp]{lem:sub:R}.
\end{prob}

\begin{lem}\ollabel{lem:sub:left-alpha}
  إذا تحقق~$P\aeq Q$ وكان كل من~$\Subst{P}{R}{y}$ و
  $\Subst{Q}{R}{y}$ معرّفًا، لزم
  $\Subst{P}{R}{y}\aeq\Subst{Q}{R}{y}$.
\end{lem}

\begin{proof}
  نقدم حقيقة مساعدة يثبتها الاستقراء القوي في عدد البانيات: لكل حد~$T$
  ومجموعة منتهية~$F$ من أسماء المتغيرات يوجد حد~$T_F\aeq T$ لا يسمى أي
  تجريد فيه باسم من~$F$. فأمر المتغير والتطبيق بيّن. وفي حالة
  $\lambd[x][U]$ نختار~$w\notin F\cup\FV{U}$ مع~$w\ne x$، ثم نستعمل
  فرض الاستقراء في~$U$ بعد أن نزيد~$\{x,w\}$ في مجموعة المنع؛ فيغدو
  إحلال~$w$ مكان~$x$ في المتن الناتج معرّفًا، وتصح إعادة تسمية الرابط
  بالقاعدة~\olref{defn:aconvone4}. ولا دور للّمة الحاضرة في هذا
  الاستقراء.

  ثم يكفي أن نعالج خطوة~$P\aconvone Q$ واحدة، ونتعدى منها إلى السلسلة.
  أما حالات التوافق الثلاث فتتبع من فرض الاستقراء ومن تعريف الإحلال.
  وفي تغيير الرابط عند الجذر نكتب
  $P=\lambd[x][N]$ و$Q=\lambd[z][\Subst{N}{z}{x}]$.
  نطبق الحقيقة المساعدة على~$R$ ومجموعة المنع~$\{x,z\}$، فنأخذ
  $R^\ast\aeq R$ ليس في روابطه اسم~$x$ ولا~$z$.
  وتفيد~\olref{lem:sub:R} أن الإحلالين يظلان معرّفين إذا جعلنا
  $R^\ast$ مكان~$R$، وأن نتيجتيهما تكافئان بحسب~$\alpha$ نتيجتيهما
  الأوليين.

  ومن شروط خطوة تغيير الرابط وكون الإحلالين معرّفين أن~$x,y,z$ مختلفة
  اثنين اثنين، وأن~$z\notin\FV{N}$ و$x,z\notin\FV{R^\ast}$.
  ويعطي الاستقراء البنيوي على~$N$، مع استعمال شروط تعريف الإحلال في
  حالة التجريد، أن الطرف الأيسر التالي معرّف، وأن
  \[
    \Subst{\Subst{N}{R^\ast}{y}}{z}{x}
    =\Subst{\Subst{N}{z}{x}}{R^\ast}{y}.
  \]
  ففي حالة~$N=y$ نعبر روابط~$R^\ast$؛ ولا يسمى شيء منها~$x$ أو~$z$.
  وسائر حالات المتغيرات ظاهرة، وأما التطبيق والتجريد فبفرض الاستقراء.
  ولا يقع~$z$ حرًا في~$\Subst{N}{R^\ast}{y}$ كذلك. فتجيز
  \olref{defn:aconvone4} أن نكتب
  \[
    \Subst{P}{R^\ast}{y}
    =\lambd[x][\Subst{N}{R^\ast}{y}]
    \aconvone
    \lambd[z][\Subst{\Subst{N}{z}{x}}{R^\ast}{y}]
    =\Subst{Q}{R^\ast}{y}.
  \]
  ثم نرجع من~$R^\ast$ إلى~$R$ بــ\olref{lem:sub:R}، ونستعمل تعدي
  $\aeq$. وتتم سلسلة التحويل بتكرار حجة الخطوة الواحدة.
\end{proof}

سبق في~\olref[sub]{sec} أن الإحلال قد لا يكون معرّفًا.
غير أن تحويل~$\alpha$ يتيح تغيير أسماء المتغيرات المقيدة حتى يصير الإحلال معرّفًا في كل حالة.
ويظل التكافؤ بحسب~$\alpha$ محفوظًا في النتيجة، كما تقرره المبرهنة التالية:

\begin{thm}\ollabel{thm:sub}
  أيًا تكن~$M$ و$R$ و$y$، يوجد~$M'$ يحقق~$M \aeq M'$ ويكون
  $\Subst{M'}{R}{y}$ معرّفًا. وإذا كان~$M'',R''$ زوجًا آخر يحقق
  $M'' \aeq M$ و$R'' \aeq R$، وكان~$\Subst{M''}{R''}{y}$ معرّفًا،
  لزم~$\Subst{M'}{R}{y} \aeq \Subst{M''}{R''}{y}$ أيضًا.
\end{thm}

\begin{proof}
  نستقرئ استقراءً قويًا في عدد بانيات الحد~$M$:
  \begin{enumerate}
  \item إن كان~$M$ متغيرًا~$z$، فالإثبات تمرين.
  \item ليكن~$M$ على صورة~$\lambd[x][N]$. نختار متغيرًا~$z$
    غير~$x$ و$y$ بحيث~$z \notin \FV{N}$ و$z \notin \FV{R}$.
    ونستعمل فرض الاستقراء أولًا في~$N$ مع الإحلال
    $\Subst{N}{z}{x}$، فنأخذ~$N_0 \aeq N$ يكون
    $\Subst{N_0}{z}{x}$ معرّفًا. ولنضع~$P=\Subst{N_0}{z}{x}$.
    حينئذ
    $\lambd[x][N] \aeq \lambd[x][N_0] \aeq \lambd[z][P]$؛
    لأن~$z\ne x$، ولأن~$z\notin\FV{N_0}=\FV{N}$، ولأن~$P$ معرّف.
    وإذ كان تحويل~$\alpha$ وإحلال متغير مكان متغير لا يبدلان عدد
    البانيات، كان~$P$ دون~$M$؛ لذلك نستعمل فرض الاستقراء القوي مرة
    أخرى في~$P,R,y$، فنأخذ
    $P'\aeq P$ يكون~$\Subst{P'}{R}{y}$ معرّفًا.
    ويرفع توافق~$\aeq$ هذا إلى
    $\lambd[z][P]\aeq\lambd[z][P']$.
    فإذا جعلنا~$M'=\lambd[z][P']$ تحقق~$M\aeq M'$، وصار
    $\Subst{M'}{R}{y}$ معرّفًا؛ إذ إحلال المتن معرّف، و$z\ne y$،
    و$z\notin\FV{R}$.

    \item وإن كان~$M$ على صورة~$(PQ)$، فالإثبات تمرين.
  \end{enumerate}

  ويبقى الانفراد. ليكن~$M''\aeq M$ و$R''\aeq R$، وليكن
  $\Subst{M''}{R''}{y}$ معرّفًا. لما كان~$M'\aeq M$، أفاد تناظر
  $\aeq$ وتعديه~$M'\aeq M''$. ومن \olref{lem:sub:R}، مطبقةً على
  $R''\aeq R$، يكون~$\Subst{M''}{R}{y}$ معرّفًا، ويتحقق
  \[
    \Subst{M''}{R''}{y}\aeq\Subst{M''}{R}{y}.
  \]
  ثم تفيد \olref{lem:sub:left-alpha}
  $\Subst{M'}{R}{y}\aeq\Subst{M''}{R}{y}$؛ وبالتعدي يلزم
  $\Subst{M'}{R}{y}\aeq\Subst{M''}{R''}{y}$.
  % تصحيح مصدر (0362-I2): استُعمل الاستقراء القوي مرتين لضمان تعريف الإحلال، وأثبت توافق الإحلال مع ألفا بعد تنضير أسماء الروابط، واستُبدلت مساواة الممثلين بتكافؤ ألفا.
\end{proof}

\begin{prob}
  أتمّ برهان~\olref[lam][syn][alp]{thm:sub}.
\end{prob}

\begin{cor}\ollabel{cor:sub}
  أيًا تكن~$M$ و$R$ و$y$، يوجد زوج~$M'$ و$R'$ يحقق~$M' \aeq M$ و$R \aeq
  R'$،
  ويكون~$\Subst{M'}{R'}{y}$ معرّفًا. ومتى وجد زوج آخر يحقق~$M'' \aeq M$ و$R'' \aeq R$
  وكان~$\Subst{M''}{R''}{y}$ معرّفًا، لزم~$\Subst{M'}{R'}{y} \aeq \Subst{M''}{R''}{y}$.
\end{cor}
% تصحيح مصدر (0362-I3): أُثبت شرط تكافؤ حدي الإحلال في الزوجين.

\begin{proof}
  يلزم هذا مباشرة من~\olref{thm:sub}.
\end{proof}

\end{document}
